\scanpage{289}
Soit $Y=\operatorname{Spec}(V)$, $V$ anneau de valuation discrète

$X$ schéma abélien sur $Y$

$y_0,y_1$ pt fermé resp. gén. de $Y$

$X_0,X_1$ les deux fibres de $X$ \ldots

$F,G$ deux sous-groupes de $X$, \uline{finis} sur $Y$,
\[
\begin{array}{ll}
\unsure{a)|i)}&F_1\not\subset G_1\deleted{\unsure{\ne 0}}\\
\unsure{a')|ii)}&F_0\subset G_0\ne 0\\
\deleted{\unsure{a'')}}\inserted{\unsure{iii)|a'')}}&pF_{\unsure{1|0}}=G_1\qquad(p=\operatorname{car}k(V))
\end{array}
\]
\[
\begin{array}{l}A=X/F\\B=X/G\end{array}
\]
\diagram{\begin{tikzcd}[row sep=small]X\arrow[r,"\alpha"]&A\\X\arrow[r,"\beta"]&B\end{tikzcd}}
hom can.

Considérons $\beta:X\longrightarrow X/G$ (canonique) alors $p\beta$ s'annule sur \deleted{$pX$, donc sur} $F$ \inserted{en vertu de (iii)}, donc se factorise par \uline{$u$}: $X/F-X/G$
\diagram{\begin{tikzcd}[column sep=small]
&X\arrow[dl,"\alpha"']\arrow[dr,"p\beta"]&\\
A\arrow[rr,"u"]&&B
\end{tikzcd}}
\deleted{Cependant} $u_1\notin p\operatorname{Hom}(A_1,B_1)$, en \deleted{\illegible{1}} se plaçant dans $\operatorname{Hom}(X_1,X_1)\otimes\mathbb Q\simeq\operatorname{Hom}(A_1,B_1)\otimes\mathbb Q$, si on avait $u=pv$, on aurait $v=\delta$, or $\delta$ ne provient pas d'un élément de $\operatorname{Hom}(A_1,B_1)$, puisque $F_1\not\subset G_1$ \deleted{[\unsure{sinon}, par (i), on aurait $F_1=0$ d'où $F_0=0$ qui contredit (ii)]}. Cependant $u_0\in p\operatorname{Hom}(A_0,\unsure{F_0|B_0})$, puisque $F_0\subset G_0$.

Les conditions dites (i)(ii)(iii) peuvent être réalisées, en égales ou inégales car, avec pour $X$ une \uline{courbe elliptique}, on prend $X_0$ sans pt d'ordre $p$, $X_1$ avec pt d'ordre $p$, et $F_1$ et $G_1$ deux ss-groupes de $X_1$, \deleted{annulés par $p$}\inserted{de \deleted{\unsure{ordre}} rang $p$ sur $k(y_1)$ et} distincts. Comme il n'y a
\figure{\begin{tikzpicture}\draw (0,1.1)--(0,1.3)--(1.9,1.45)--(1.92,1.55);\draw (0.1,0)--(0.1,0.33)--(2.6,0.31)--(2.7,0.1);\end{tikzpicture}}
\mspage{288}
